% page 486 du fac-simile (image p-485 du scan)
% bloc 195.6 x 250.6 mm ; marge gauche 3.0 mm
\pgc{-17.780mm}{0.000mm}{
\l{}
\l{}
\l{}
\l{}
\l{}
\l{}
\l{}
\l{\cel{17}\sou{quinto}. I. (\sou{muziko}).Intervalo de kin noti intersequanta,}
\l{\cel{20}la du extrema inkluzite. - Grado kinesma di la skalo dia-}
\l{\cel{20}tonika. - II. (\sou{skerm}.)Pareo simpla, qua korespondas a la}
\l{\cel{20}engajo (adinterne) di la lineo alta. - DEFIRS.}
\l{}
\l{\cel{17}\sou{quintalo}. (\sou{mezur-sistemo metrala}). Maso 100-kilograma. -EFIS}
\l{}
\l{\cel{17}\sou{quintana}. (\sou{patol.}) Febro qua riaparas singla-kindie. -}
\l{}
\l{\cel{17}\sou{quintesenco}. I. (a) (\sou{olim)}.\cel{2}Substanco eterala, egardatakom}
\l{\cel{20}la elemento kinesma e maxim subtila. (b). (\sou{nun)}. La parto}
\l{\cel{20}maxim subtila di substanco. - II. (\sou{metaf}.)Lo maxim bona,}
\l{\cel{20}maxim precoza, maxim perfekta. - DEFIRS.}
\l{}
\l{\cel{17}\sou{quinteto}. Muzikajo de kin parti koncertala. - DEFIS.}
\l{}
\l{\cel{17}\sou{quintiko}. (\sou{matem.}) Kurvo di la klaso kinesma. EFIS.}
\l{}
\l{\cel{17}\sou{quintiliono}. (\sou{matem.)}\cel{3}Mil quadrilioni. (1030).}
\l{}
\l{\cel{17}\sou{quiproquo}. Eroro konseque de qua onu egardas kozo o perso-}
\l{\cel{20}no kom altra. - DEFIS.}
\l{}
\l{\cel{17}\sou{quirito}. (\sou{che la Romani, antique}). Civitano Romana qua rezi-}
\l{\cel{20}dis en Roma, kontraste kun ti qui esis servanta en la}
\l{\cel{20}armei. - DEFIRS.}
\l{}
\l{\cel{17}quirlar. (trans.) Batar (liquido, ovi, kremo) per vergo por}
\l{\cel{20}mixar, igar augmenter volumine, spumifar, e c.) - DE.}
\l{}
\l{\cel{17}\sou{qui}ta. Tote liberigata de obligeso pekuniala od etikala.}
\l{\cel{20}- DEFS.}
\l{}
\l{\cel{17}\sou{quo.}\cel{2}Neutra formo di \textquotedbl{}qua\textquotedbl{}.}
\l{}
\l{\cel{17}\sou{quociento}. (\sou{aritm.}) Rezultajo de la divido da nombro per}
\l{\cel{20}altra. - DEFIS.}
\l{}
\l{\cel{17}\sou{quodlibeto}. Dico triviala, mokachanta, joko mala-mora, mala}
\l{\cel{20}vorto-ludo. - F.}
\l{}
\l{\cel{17}+\sou{quoniam}. Ta vorto indikas la motivo, la kauzo, la justifi-}
\l{\cel{20}kivo ; de la instanto kande agnoskas kom valida ke...- L.}
\l{}
\l{\cel{17}\sou{quoto}. La parto di singlu en la repartiso di pekunio-quanto}
\l{\cel{20}pagenda o recevenda da ensemblo de personi. -\cel{2}DEFIS.}
\l{}
\l{}
\l{}
\l{}
\l{\cel{48}---}
}
